% concatenated from 240210221.tex
%%%%%%%%%%%%%%%%%%%%%%% file template.tex %%%%%%%%%%%%%%%%%%%%%%%%%
%
% This is a general template file for the LaTeX package SVJour3
% for Springer journals.          Springer Heidelberg 2010/09/16
%
% Copy it to a new file with a new name and use it as the basis
% for your article. Delete % signs as needed.
%
% This template includes a few options for different layouts and
% content for various journals. Please consult a previous issue of
% your journal as needed.
%
%%%%%%%%%%%%%%%%%%%%%%%%%%%%%%%%%%%%%%%%%%%%%%%%%%%%%%%%%%%%%%%%%%%
%
% First comes an example EPS file -- just ignore it and
% proceed on the \documentclass line
% your LaTeX will extract the file if required
%\begin{filecontents*}{example.eps}
%!PS-Adobe-3.0 EPSF-3.0
%%BoundingBox: 19 19 221 221
%%CreationDate: Mon Sep 29 1997
%%Creator: programmed by hand (JK)
%%EndComments
%gsave
%newpath
%  20 20 moveto
%  20 220 lineto
%  220 220 lineto
%  220 20 lineto
%closepath
%2 setlinewidth
%gsave
%  .4 setgray fill
%grestore
%stroke
%grestore
%\end{filecontents*}
%
\RequirePackage{fix-cm}
%
%\documentclass{svjour3}                     % onecolumn (standard format)
%\documentclass[smallcondensed]{svjour3}     % onecolumn (ditto)
\documentclass[smallextended]{svjour3}       % onecolumn (second format)
%\documentclass[twocolumn]{svjour3}          % twocolumn
%
\smartqed  % flush right qed marks, e.g. at end of proof
%
\usepackage{graphicx}
\usepackage {color}
\usepackage {tikz}


%
% \usepackage{mathptmx}      % use Times fonts if available on your TeX system
%
% insert here the call for the packages your document requires
%\usepackage{latexsym}
% etc.
%
% please place your own definitions here and don't use \def but
% \newcommand{}{}
%
% Insert the name of "your journal" with
% \journalname{myjournal}
%
\smartqed  % flush right qed marks, e.g. at end of proof
%
\usepackage{graphicx}
\usepackage{amsmath}
\usepackage{amssymb}
\usepackage{mathtools}
\def\ST{\songti\rm\relax}
\def\R{{\mathbb R}}
\def\Z{{\mathbb Z}}
\def\ST{{\rm s.t.}}
\def\argmin{{\rm argmin}}



\newtheorem{assumption}{Assumption}
%\newtheorem{example}{Example}
\newtheorem{observation}{Observation}
\begin{document}

\title{Convergence Rate of Projected Subgradient Method with Time-varying Step-sizes}
%Subgradient Method in Nonsmooth Convex Optimization: A Weak Ergodic Convergence Rate}
%about the article that should go on the front page should be
%placed here. General acknowledgments should be placed at the end of the article.}

%\subtitle{Do you have a subtitle?\\ If so, write it here}

\titlerunning{On Step-size of Subgradient Method}        % if too long for running head

\author{Zhihan Zhu \and Yanhao Zhang \and Yong Xia %etc.
}
%\authorrunning{Short form of author list} % if too long for running head

\institute{Zhihan Zhu \and Yanhao Zhang \and Yong Xia (corresponding author) \at School of Mathematical Sciences, Beihang University, Beijing 100191, People's Republic of China.
	%Tel.: +123-45-678910\\
	%Fax: +123-45-678910\\
	\email{\{zhihanzhu,\ yanhaozhang,\ yxia\}@buaa.edu.cn}
	  }     %  \\
%             \emph{Present address:} of F. Author  %  if needed
%           \and
%           S. Author \at
%              LMIB of the Ministry of Education, School of Mathematical Sciences, Beihang University, Beijing 100191, People's Republic of China
%               \email{yxia@buaa.edu.cn}


\date{Received: date / Accepted: date}
% The correct dates will be entered by the editor
\maketitle

\begin{abstract}
We establish the optimal ergodic convergence rate for the classical projected subgradient method with  time-varying step-sizes. This convergence rate remains the same even if we slightly increase the weight of the most recent points, thereby relaxing the ergodic sense.
	\keywords{Subgradient method \and step-size \and ergodic convergence rate \and nonsmooth convex optimization}
	% \PACS{PACS code1 \and PACS code2 \and more}
	\subclass{90C25, 90C30}
\end{abstract}

\section{Introduction}
Consider the nonsmooth convex optimization problem
\begin{equation*}
	x^*\in\argmin_{x\in\mathcal{X}}f(x),
\end{equation*}
where $\mathcal{X}\subset\R^n$ is a compact convex set that is contained within the Euclidean ball $B(x^*,R)$, and $f$ is (possibly nonsmooth) convex and Lipschitz on $\mathcal{X}$, i.e., there is an $L>0$ such that $\Vert g\Vert\le L$ for any $g\in\partial f(x)\neq\emptyset$ and $x\in \mathcal{X}$.
%\begin{equation*}
%	{\rm \Pi}_{\mathcal{X}}(y)=,
%\end{equation*}

The classical projected subgradient method (PSG) iterates the following equations for $t\ge 1$:
\begin{equation*}
\left\{	\begin{aligned}
		y_{t+1}&=x_t-\eta_t g_t, ~\text{where } g_t\in\partial f(x_t),\\
		x_{t+1}&=  \argmin_{x\in\mathcal{X}}\Vert x-y_{t+1} \Vert.	\end{aligned} \right.
\end{equation*}
Utilizing  the following constant step size
	\begin{equation}
		\eta_s\equiv\frac{R}{L\sqrt{t}}, ~s=1,\cdots,t,  \label{size1}
	\end{equation}	
PSG achieves an optimal ergodic convergence rate (see, for example, \cite{nesterov2018lectures,bubeck2015convex,MITlecture}) expressed as \footnote{ $f ((\sum_{s=1}^t x_s)/t)$ can be replaced with $\min_{s=1,\cdots,t}f(x_s)$ based on similar analysis.}
\begin{equation}
		f\left(\frac{\sum_{s=1}^t x_s}{t}\right)-f(x^*)\le \frac{RL}{\sqrt{t}}. \label{rate1}
\end{equation}

A more practical approach is to use a step-size that monotonically decreases towards $0$.
One such time-varying step-size, motivated by \eqref{size1} and suggested in \cite{nesterov2018lectures,bubeck2015convex}, is expressed as follows:
\begin{equation}
	\eta_s=\frac{R}{L\sqrt{s}},~s=1,\cdots,t.\label{size2}
\end{equation}
However, using this step size results in sub-optimal ergodic convergence rate \cite{nesterov2018lectures,bubeck2015convex}, as it adds an additional $\log(t)$ factor compared to the right-hand side of \eqref{rate1}.
%Note that computing the weighted average of the iterates in the second half of the sequence could yield the optimal convergence rate \cite[Corollary 3.2]{Lan}.


The contribution of this note is to demonstrate that PSG with the time-varying step-size \eqref{size2} indeed achieves the following optimal convergence rate.
\begin{theorem} \label{main}
PSG with the time-varying step-size \eqref{size2} satisfies
 \begin{equation*}
	f\left(\frac{\sum_{s=1}^t x_s}{t}\right)-f(x^*)\le \frac{3RL}{2\sqrt{t}}.
\end{equation*}
 \end{theorem}
In Section 2, we present a more generalized convergence analysis, allowing us to provide some insightful observations.

%Since the subgradient direction does not guarantee decrease of the function value, the choice of step size becomes particularly important in the subgradient method. Constant step size scheme requires pre-determining the maximum number of iterations based on the desired accuracy, which can be cumbersome in practice. Time-varying step sizes appear more flexible. However, a somewhat counterintuitive conclusion is that, when the objective function is convex and Lipschitz, the subgradient method with a constant step size is superior to the time-varying step size subgradient method \cite{nesterov2018lectures,bubeck2015convex}, with the convergence rate difference of a logarithmic factor. Although \cite{MITlecture} provides a solution to eliminate the logarithmic factor in the convergence rate of time-varying step size subgradient methods, the approach is intricate and not essential.
%
%In this paper, we use novel techniques to prove the convergence rate of  subgradient methods with time-varying step size. We demonstrate that, through concise operations, subgradient methods with time-varying step size can achieve optimal convergence rates. This indicates that time-varying step size scheme is also optimal and easier to implement in practice. We introduce the concept of weak ergodic convergence and explain why the original scheme [2,3] fail to unleash the potential of time-varying step size methods. The conclusions can be extended to techniques such as mirror descent, dual averaging, and other nonsmooth convex optimization methods with time-varying step size.

%\section{Weak Ergodic Convergence Rate of Subgradient Method}
\section{Analysis}
%In this section, we investigate the weak ergodic convergence rate of subgradient method in nonsmooth convex optimization. Consider the optimization problem
%\begin{equation*}
%	\min_{x\in\mathcal{X}}f(x),
%\end{equation*}
%where $\mathcal{X}\subset\R^n$ is a compact convex set, and $f$ is such that for any $x\in\mathcal{X}$, $\partial f(x)\neq\emptyset$ and Lipschitz on $\mathcal{X}$. We denote $x^*\in\argmin_{x\in\mathcal{X}}f(x)$  and assume that $\mathcal{X}$ is contained in the Euclidean ball $B(x^*,R)$. We define the projection operator $\rm \Pi_{\mathcal{X}}$ on $\mathcal{X}$ by
%%\begin{equation*}
%%	{\rm \Pi}_{\mathcal{X}}(y)=,
%%\end{equation*}
%and thereby obtain the projection subgradient method
%\begin{equation*}
%	\begin{aligned}
%		y_{t+1}&=x_t-\eta_t g_t, \text{where } g_t\in\partial f(x)\\
%		x_{t+1}&=  \argmin_{x\in\mathcal{X}}\Vert x-y_{t+1} \Vert.
%	\end{aligned}
%\end{equation*}
Consider PSG with a general step-size $\eta_s$, which is
assumed to be positive and non-increasing. We can establish the following convergence rate.
\begin{theorem}\label{Convergence}
For any fixed $k\ge -1$,
PSG with a positive and non-increasing step-size sequence $\{\eta_s\}$ satisfies
	\begin{equation}
		f\left(\frac{\sum_{s=1}^t \frac{1}{\eta_s^k}x_s}{\sum_{s=1}^t \frac{1}{\eta_s^k}}\right)-f(x^*)\le \frac{\frac{R^2}{\eta_t^{k+1}}+ L^2 \sum_{s=1}^t\frac{1}{\eta_s^{k-1}}}{2\sum_{s=1}^t\frac{1}{\eta_s^k}}. \label{bound}
	\end{equation}
\end{theorem}	
\begin{proof}
	According to the definition of subgradient, we have
	\begin{eqnarray}
			f(x_s)-f(x^*)&\le& g_s^T(x_s-x^*) \nonumber \\
			&=& \frac{1}{\eta_s}(x_s-y_{s+1})^T(x_s-x^*) \nonumber\\
			&=& \frac{1}{2\eta_s}(\Vert x_s-y_{s+1} \Vert^2+\Vert x_s-x^* \Vert^2-\Vert y_{s+1}-x^* \Vert^2)\label{eq1}\\
			&=& \frac{1}{2\eta_s}(\Vert x_s-x^* \Vert^2-\Vert y_{s+1}-x^* \Vert^2)+\frac{\eta_s}{2}\Vert g_s \Vert^2 \nonumber\\
&\le& \frac{1}{2\eta_s}(\Vert x_s-x^* \Vert^2-\Vert x_{s+1}-x^* \Vert^2)+\frac{\eta_s}{2}L^2,\label{ineq1}
	\end{eqnarray}
where \eqref{eq1} follows from the elementary identity $2a^Tb=\Vert a\Vert^2+\Vert b\Vert^2-\Vert a-b\Vert^2$, and \eqref{ineq1} holds since $\Vert g_s \Vert \le L$ and
	\begin{equation*}
		\Vert y_{s+1}-x^* \Vert^2\ge \Vert x_{s+1}-x^* \Vert^2,
	\end{equation*}
which is implied by  the projection theorem.
%	Since $f$ is $L$-Lipschitz on $\mathcal{X}$, $\Vert g_s \Vert \le L$, we have
%	\begin{equation*}
%		f(x_s)-f(x^*)\le \frac{1}{2\eta_s}(\Vert x_s-x^* \Vert^2-\Vert x_{s+1}-x^* \Vert^2)+\frac{\eta_s}{2}L^2.
%	\end{equation*}
%Consequently, for any $k\ge -1$,
%	\begin{equation*}
%		\frac{1}{\eta_s^k}(f(x_s)-f(x^*))\le \frac{1}{2\eta_s^{k+1}}(\Vert x_s-x^* \Vert^2-\Vert x_{s+1}-x^* \Vert^2)+\frac{1}{2\eta_s^{k-1}}L^2.
%	\end{equation*}
%	Summing up the inequality over $s$, we get

Consequently, we have
	\begin{eqnarray}
&&\left(\sum_{s=1}^t\frac{1}{\eta_s^k}\right)\left(		f\left(\frac{\sum_{s=1}^t \frac{1}{\eta_s^k}x_s}{\sum_{s=1}^t \frac{1}{\eta_s^k}}\right)-f(x^*)\right) \nonumber \\
&\le&	\sum_{s=1}^t\frac{1}{\eta_s^k}(f(x_s)-f(x^*))  ~~~~({\rm since~}f{\rm~is~convex})\nonumber\\
&\le& \sum_{s=1}^t\frac{1}{2\eta_s^{k+1}}(\Vert x_s-x^* \Vert^2-\Vert x_{s+1}-x^* \Vert^2)+\sum_{s=1}^t\frac{1}{2\eta_s^{k-1}}L^2 \label{ineq2}\\
			&=& \frac{1}{2\eta_1^{k+1}}\Vert x_1-x^* \Vert^2+\sum_{s=2}^t(\frac{1}{2\eta_s^{k+1}}-\frac{1}{2\eta_{s-1}^{k+1}})\Vert x_s-x^* \Vert^2 \nonumber\\
&&-\frac{1}{2\eta_{t}^{k+1}}\Vert x_{t+1}-x^* \Vert^2+\sum_{s=1}^t\frac{1}{2\eta_s^{k-1}}L^2\nonumber\\
			&\le& \frac{R^2}{2\eta_1^{k+1}}+\frac{R^2}{2}\sum_{s=2}^t(\frac{1}{\eta_s^{k+1}}-\frac{1}{\eta_{s-1}^{k+1}})+
\sum_{s=1}^t\frac{1}{2\eta_s^{k-1}}L^2\label{ineq3}\\
			&=& \frac{R^2}{2\eta_t^{k+1}}+\sum_{s=1}^t\frac{1}{2\eta_s^{k-1}}L^2,\nonumber
		 \end{eqnarray}
%	Since $f$ is convex on $\mathcal{X}$,
%	\begin{equation*}
%(\sum_{s=1}^t\frac{1}{\eta_s^k})(		f(\frac{\sum_{s=1}^t \frac{1}{\eta_s^k}x_s}{\sum_{s=1}^t \frac{1}{\eta_s^k}})-f(x^*))\le \frac{\frac{R^2}{2\eta_t^{k+1}}+\frac{L^2}{2}\sum_{s=1}^t\frac{1}{\eta_s^{k-1}}}{\sum_{s=1}^t\frac{1}{\eta_s^k}}.
%	\end{equation*}
where \eqref{ineq2} follows from \eqref{ineq1}, and \eqref{ineq3} holds since $1/\eta_s^{k+1}-1/\eta_{s-1}^{k+1}\ge0$ when $k\ge -1$. The proof is complete.  \qed
\end{proof}	

\begin{remark}\label{r1}
By setting $k=-1$ in Theorem \ref{Convergence}, the upper bound on the right-hand side \eqref{bound} simplifies to
	\begin{equation*}
\frac{R^2+L^2\sum_{s=1}^t\eta_s^2}{2\sum_{s=1}^t\eta_s},
	\end{equation*}
	which is exactly the same as the result presented in \cite{nesterov2018lectures}. Then PSG with the time-varying step-size \eqref{size2} satisfies
\begin{equation}
	f\left(\frac{\sum_{s=1}^t \frac{1}{ \sqrt{s}}x_s}{\sum_{s=1}^t \frac{1}{\sqrt{s}}}\right)-f(x^*)\le \frac{2RL+RL\log t}{4(\sqrt{t+1}-1)}, \label{bound2}
\end{equation}
which is sub-optimal. Note that computing the weighted average of the iterates in the second half of the sequence yields the optimal convergence rate \cite[Corollary 3.2]{Lan}.
\end{remark}

%\begin{remark}\label{r2}
%	In Theorem \ref{Convergence}, if we consider the constant step size
%	\begin{equation*}
%		\eta_s=\frac{R}{L\sqrt{t}}, s=1,...,t,
%	\end{equation*}
%	then we have (see also \cite{nesterov2018lectures,bubeck2015convex,MITlecture})
%	\begin{equation*}
%		f(\frac{\sum_{s=1}^t x_s}{t})-f(x^*)\le \frac{RL}{\sqrt{t}}.
%	\end{equation*}
%	
%\end{remark}

\begin{remark}\label{r3}
By setting $k=0$ in Theorem \ref{Convergence},  we can immediately obtain the optimal convergence rate, as presented in Theorem \ref{main}.
\end{remark}

%by setting $k=-1$, which is also the result obtained in \cite{nesterov2018lectures,bubeck2015convex,MITlecture}. Hence, \cite{nesterov2018lectures} claims that the subgradient method with time-varying step size is sub-optimal, since there is an extra $\log t$ factor. \cite{MITlecture} proposed a method by ignoring the first $t/2$ point and summing from $t/2$ to $t$ to eliminate the $\log t$ factor. However, by simply setting $k=0$ in Theorem \ref{Convergence}, we will see \begin{equation*}
%	f(\frac{\sum_{s=1}^t x_s}{t})-f(x^*)\le \frac{3RL}{2\sqrt{t}}.
%\end{equation*}
%This implies that for the case of time-varying step size, we can achieve the same optimal convergence rate as with constant step size by simply calculating $\frac{\sum_{s=1}^t x_s}{t}$, without the need for any additional operations. Why could this happen?
%
%We notice that the original weighting scheme $\frac{\sum_{s=1}^t \eta_sx_s}{\sum_{s=1}^t \eta_s}$ assigns larger weights to the initial points and smaller weights to the most recent points. However, in the optimization process, the value of the most recent points must be higher than that of the initial points. Therefore, the original weighting scheme cannot fully unleash the potential of the time-varying step size approach. By taking the average $\frac{\sum_{s=1}^t x_s}{t}$, we effectively increase the weight of the most recent points, resulting in a better convergence rate. The method proposed in \cite{MITlecture}, in fact, also artificially ignores the first $t/2$ points, assigning weights to the latter $t/2$ points, achieving a similar effect as increasing k.
%
%We refer to the convergence rate achieved by assigning higher weights to the most recent points as weak ergodic convergence rate. When weight is only assigned to the latest point, the weak ergodic convergence rate degenerates into non-ergodic convergence rate.

\begin{remark}
By setting any $k$ such that $k>-1$ in Theorem \ref{Convergence},  the convergence rate of  $f((\sum_{s=1}^t  x_s/\eta_s^k)/(\sum_{s=1}^t 1/ \eta_s^k))-f(x^*)$ will be $\mathcal{O}(1/\sqrt{t})$, without the presence of a $\log t$ factor. In comparison with the sub-optimal case \eqref{bound2}, when $k>0$,  the
 weighting scheme $(\sum_{s=1}^t  x_s/\eta_s^k)/(\sum_{s=1}^t 1/ \eta_s^k)$ assigns smaller weights to the initial points and larger weights to the most recent points. This new result can be referred to as the ``weak'' ergodic convergence rate.
\end{remark}

\begin{remark}\label{r4}
	We can apply the same proof techniques to extend the conclusion of weak ergodic convergence to mirror descent, Nesterov's dual averaging, and other schemes with time-varying step sizes for solving nonsmooth convex optimization, see \cite{nesterov2018lectures,bubeck2015convex}.

\end{remark}

%\section{Conclusion}
%In this paper, we employ novel proof techniques to establish the convergence rate of the subgradient method with time-varying step sizes. By directly computing the average, we can achieve optimal convergence rate same as constant step size schemes, without the need for any additional operations. This is because the original weighted averaging scheme assigns excessively high weights to the initial points and smaller weights to the most recent points, resulting in a convergence rate difference of a logarithmic factor compared to constant step size schemes and failing to unleash the potential of time-varying step size approaches. By slightly increasing the weight of the most recent points, we can attain the optimal convergence rate. We refer to this method, which improves convergence rate by elevating the weight of the most recent points, as weak ergodic convergence rate. When weight is only assigned to the latest point, weak ergodic convergence rate degenerates into non-ergodic convergence rate. The concept of weak ergodic convergence can also be applied to enhance mirror descent, dual averaging, and other methods with time-varying step size in nonsmooth optimization.

\section*{Funding}
This research was supported by National Key R\&D Program of China under grant 2021YFA1003300, the National Natural Science Foundation of China under grant 12171021, and the Fundamental Research Funds for the Central Universities.


\section*{Data Availability}
The manuscript has no associated data.

\begin{thebibliography}{9}
%\bibliographystyle{spmpsci}      % mathematics and physical sciences
\bibitem{nesterov2018lectures}
Y. Nesterov, Lectures on convex optimization, volume 137, Springer, 2018.
\bibitem{bubeck2015convex}
S. Bubeck, Convex optimization: Algorithms and complexity. Foundations
and Trends Trends{\textregistered} in Machine Learning, 8(3-4): 231-357, 2015.
\bibitem{MITlecture}
P. Rigollet, K. Li, Convex optimization for machine learning. Massachusetts
Institute of Technology: MIT OpenCouseWare, https://ocw.mit.edu/, Oct. 2015.
\bibitem{Lan}
G. Lan, First-order and Stochastic Optimization Methods for Machine Learning, Springer-Nature, 2020
\end{thebibliography}
%\bibliography{references}

\end{document}
% end of file template.tex

